\section{CNN}
\subsection{CNN简介}
本章的主题是卷积神经网络（Convolutional Neural Network，CNN），
之前介绍的神经网络中，相邻层的所有神经元之间都有连接，这称为全连接网络（fully-connected Network）。

\begin{figure}[H]
	\centering
	\includegraphics[width=0.9\textwidth, keepaspectratio]{4-1.png}
	\caption{全连接网络示例}
\end{figure}

\begin{figure}[H]
	\centering
	\includegraphics[width=0.9\textwidth, keepaspectratio]{4-2.png}
	\caption{卷积神经网络示例}
\end{figure}

相较全连接网络，CNN中新出现了卷积层（Convolution层）和池化层（Pooling层），
CNN 的层的连接顺序是“Conv - ReLU - Pooling”（Pooling层有时会被省略）。
还需要注意的是，在 CNN 中，靠近输出的层中使用了之前的“Affine - ReLU”组合。
此外，最后的输出层中使用了之前的“Affine - Softmax”组合，这些都是 CNN 中比较常见的结构。

\subsection{卷积层}
全连接层（Affine层）存在什么问题呢？那就是数据的形状被“忽视”了。比如，输入数据是图像时，
图像通常是高、长、rgb通道方向上的3维形状。但是，向全连接层输入时，需要将3维数据拉平为1维数据。
图像是3维形状，这个形状中含有重要的空间信息。比如，空间上邻近的像素为相似的值，形状中可能隐藏有值得提取的特征。

卷积层可以保持形状不变。当输入数据是图像时，卷积层会以3维数据的形式接收输入数据，并同样以3维数据的形式输出至下一层。
利用卷积层，有可能正确理解输入数据的形状的信息，卷积层的输入输出数据称为特征图（feature map）。
其中，输入数据称为输入特征图（input feature map），输出数据称为输出特征图（output feature map）。

\begin{figure}[H]
	\centering
	\includegraphics[width=0.5\textwidth, keepaspectratio]{4-3.png}
	\caption{卷积运算示例}
\end{figure}

卷积层进行的处理就是卷积运算，相当于图像处理中的“滤波器运算”，滤波器也被称为卷积核，
卷积运算以一定间隔滑动滤波器的窗口并应用，这个计算也被称为乘积累加运算。

\begin{figure}[H]
    \centering
    % 第一行
    \begin{subfigure}[b]{0.45\textwidth}
        \centering
        \includegraphics[width=0.9\textwidth]{4-4a.png}
        \caption{step 1}
        \label{fig:a}
    \end{subfigure}
    \hfill
    \begin{subfigure}[b]{0.45\textwidth}
        \centering
        \includegraphics[width=0.9\textwidth]{4-4b.png}
        \caption{step 2}
        \label{fig:b}
    \end{subfigure}

    \vspace{0.5cm} % 行间距

    % 第二行
    \begin{subfigure}[b]{0.45\textwidth}
        \centering
        \includegraphics[width=0.9\textwidth]{4-4c.png}
        \caption{step 3}
        \label{fig:c}
    \end{subfigure}
    \hfill
    \begin{subfigure}[b]{0.45\textwidth}
        \centering
        \includegraphics[width=0.9\textwidth]{4-4d.png}
        \caption{step 4}
        \label{fig:d}
    \end{subfigure}

    \caption{卷积运算流程说明}
    \label{fig:grid}
\end{figure}

CNN中也存在偏置，偏置通常只有1个，这个值会被加到应用了滤波器的所有元素上。

\begin{figure}[H]
	\centering
	\includegraphics[width=0.9\textwidth, keepaspectratio]{4-5.png}
	\caption{偏置示例}
\end{figure}

在进行卷积层的处理之前，有时要向输入数据的周围填入固定的数据（比如0），这称为填充（padding），主要是为了调整输出的大小。
如下图所示，通过应用了幅度为1的填充，大小为(4, 4)的输入数据变成了(6, 6)的形状，
然后，应用大小为(3, 3)的滤波器，生成了大小为(4, 4)的输出数据。

\begin{figure}[H]
	\centering
	\includegraphics[width=0.75\textwidth, keepaspectratio]{4-6.png}
	\caption{填充示例}
\end{figure}

应用滤波器的位置间隔称为步幅（stride）。下图为一个步幅为2卷积运算示例。

\begin{figure}[H]
    \centering
    % 第一行
    \begin{subfigure}[b]{0.45\textwidth}
        \centering
        \includegraphics[width=0.9\textwidth]{4-7a.png}
        \caption{step 1}
    \end{subfigure}
    \hfill
    \begin{subfigure}[b]{0.45\textwidth}
        \centering
        \includegraphics[width=0.9\textwidth]{4-7b.png}
        \caption{step 2}
    \end{subfigure}

    \caption{步幅示例}
\end{figure}

图像是3维数据，除了高、长方向之外，还需要处理通道方向，在3维数据的卷积运算中，输入数据和滤波器的通道数要设为相同的值。
最后数据输出是1张特征图，如果要数据输出在通道方向上也拥有维度，那么就需要用到多个滤波器，生成多个特征图。

\subsection{池化层}
池化是缩小高、长方向上的空间的运算，下图是按步幅2进行2 × 2的Max池化时的处理顺序。
一般来说，池化的窗口大小会和步幅设定成相同的值。
除了Max池化之外，还有Average池化、随机池化等，在图像识别领域，主要使用Max池化。

\begin{figure}[H]
    \centering
    % 第一行
    \begin{subfigure}[b]{0.4\textwidth}
        \centering
        \includegraphics[width=0.9\textwidth]{4-8a.png}
        \caption{step 1}
    \end{subfigure}
    \hfill
    \begin{subfigure}[b]{0.4\textwidth}
        \centering
        \includegraphics[width=0.9\textwidth]{4-8b.png}
        \caption{step 2}
    \end{subfigure}

    \vspace{0.5cm} % 行间距

    % 第二行
    \begin{subfigure}[b]{0.4\textwidth}
        \centering
        \includegraphics[width=0.9\textwidth]{4-8c.png}
        \caption{step 3}
    \end{subfigure}
    \hfill
    \begin{subfigure}[b]{0.4\textwidth}
        \centering
        \includegraphics[width=0.9\textwidth]{4-8d.png}
        \caption{step 4}
    \end{subfigure}

    \caption{池化层流程说明}
\end{figure}

池化层的特征
\begin{itemize}
	\item 没有要学习的参数:池化层和卷积层不同，没有要学习的参数。
	\item 通道数不发生变化:经过池化运算，输入数据和输出数据的通道数不会发生变化。
	\item 对微小的位置变化具有鲁棒性:输入数据发生微小偏差时，池化仍会返回相同的结果。
\end{itemize}